% Shared appendix block for both venue versions of this paper.
%
% This file is the single source of truth for every derivation, appendix table and disclosed
% defect. `paper.tex` (AAS Journals) and `paper_icml.tex` (ICML format) both \input it, so the
% two versions cannot drift apart on anything below. Edit it here, not in either paper.
% It is included after the bibliography in both, so appendix cross-references resolve either way.

\appendix
\section{Differentiability of HMMP}\label{app:diff}
\subsection{Gradient flow through the discrete traversal}
The octree traversal involves discrete decisions (distance-threshold comparisons) that are
not differentiable. We decouple the geometry-only octree construction (a fixed computational
graph) from the learned message passing: the tree and interaction list are computed once
from the input positions and treated as constants during training, so gradients flow only
through the continuous message functions (the near-field $k$-NN aggregation and the
multipole $\mathrm{UPDATE}$). This is a straight-through-estimator stand-in for the discrete
open/close choices; they receive no gradient, as they are determined by geometry. We do not
claim differentiability w.r.t.\ the tree structure itself; the error introduced by the
hierarchical approximation, rather than by differentiability, is what Proposition~\ref{prop:bound} bounds.

\subsection{Framework independence}
The Barnes--Hut octree and multipole expansion are a geometry-only preprocessing module
independent of the GNN framework: it emits (i) near-field neighbour pairs and (ii) far-field
cell-to-node pairs, which are converted into standard \texttt{edge\_index} tensors usable by
PyG/DGL-style scatter-add message passing. Our experiments use the PyTorch implementation;
the octree core is framework-agnostic.

\subsection{Measured gradient error through the operator}
The two subsections above are arguments, not measurements. Section~\ref{sec:vjp} supplies the
measurement: it differentiates a density-weighted tidal energy through the exact field and through
the operator, and compares the two vector--Jacobian products. Differentiability and the
\emph{accuracy} of the gradient delivered are separate questions, and only the second is what a
reconstruction depends on. The adjoint sum \emph{contracts} the per-node tidal errors: the
amplification factor is below one in every $p{=}0$ configuration measured, with a maximum of $0.76$ at
$\theta\le0.70$, so the differentiability argued for here is delivered at an error no worse than the
tensor's, and the gradient a reconstruction descends on is \emph{more} accurate than the tensor it
came from. Because the tree is fixed and carries no gradient (above), this also bounds where the
error can originate: it is the far field's \emph{value} that is wrong, not the path the gradient
takes through the traversal.

\noindent\textbf{What an earlier version of this appendix reported, and why it changed.} It reported a
catastrophic failure at $\theta=0.70$: gradient error reaching $\infty$ and a direction cosine of
$0.000000$ on a third of the seeds, with the centre-of-mass acceptance test implicated.
\textbf{That result is withdrawn} --- it was defect D8 (Appendix~\ref{app:provenance}), not a property
of the approximation. With the geometry fixed the same
code path gives $0$ failures in $15$ seeds at $\theta=0.70$, an amplification of at most $0.76$ and
$\cos\ge0.99999$ --- which is why Table~\ref{tab:vjp} carries a $\theta=0.70$ row at all. The
diagnostic worth keeping is that the measurement which exposed it varied a parameter the theory says
should barely matter.

\noindent\textbf{Scope, and a defect found on first execution.} Every E3 number is on synthetic
lognormal clouds, not the CAMELS snapshot. The cloud is clustered with density-proportional masses
($\mathrm{std}=0.657$), whereas the real subsample is equal-mass. The script had never been executed
before that measurement, and its first run revealed that the per-node denominator was computed once
at $p{=}0$ \emph{outside} the multipole-order loop and reused for $p{=}1$, so the $p{=}1$
amplification divided a $p{=}1$ gradient error by a $p{=}0$ tensor error; at $n=10^{5}$,
$\theta=0.5$ this read $0.389$ where the order-matched ratio is $0.651$. The defect understated the
amplification and is fixed; the before-state artifact is kept. The measurement also required adding
\texttt{-{}-opening-mode} to the script.

\section{Proof of Proposition~\Propref}
\label{app:proof}

\subsection{Setup and notation}
Fix a graph $G$ with node features $\{h_i\}$. For each node $i$, the interaction list is
partitioned into a near set $\mathcal{N}(i)$ (exact leaf sums and $k$-NN edges) and a far set
$\mathcal{C}_i$ of cells. A cell $c$ with size $R_c$ at distance $r_{ic}$ is placed in
$\mathcal{C}_i$ iff $R_c/r_{ic}<\theta$ (the opening criterion); otherwise it is opened and
its children are considered. The physical operator evaluated by HMMP is
\begin{equation}
a_i^{\mathrm{phys}}
= \underbrace{\sum_{j \in \mathcal{N}(i)} m_{j \to i}}_{\text{exact near field}}
+ \underbrace{\sum_{c \in \mathcal{C}_i} \hat{a}_{c \to i}^{(p)}}_{\text{truncated far field}},
\end{equation}
where $\hat{a}_{c \to i}^{(p)}$ is the order-$p$ multipole approximation to the cell's
contribution and $m_{j \to i} = m_j (x_j - x_i)/\|x_j - x_i\|^3$. The exact operator replaces
$\hat{a}_{c \to i}^{(p)}$ by the exact sum over all particles in cell $c$.

\paragraph{Assumption A1 (Lipschitz update).}
The per-layer update $\phi$ is $\ell$-Lipschitz in its input: for any $u,v$,
$\|\phi(u)-\phi(v)\|\le\ell\|u-v\|$.

\paragraph{Assumption A2 (contractive aggregation).}
The aggregation operator $\mathrm{AGG}$ is $\alpha$-contractive: for any two feature sets
$\{h_j\},\{h_j'\}$,
$\|\mathrm{AGG}(\{h_j\})-\mathrm{AGG}(\{h_j'\})\|\le\alpha\max_j\|h_j-h_j'\|$.
We additionally require the \emph{effective} per-layer contraction to satisfy
$\beta:=\ell\alpha<1$; this is the condition under which the geometric series below
converges.

\paragraph{Assumption A3 (multipole remainder).}
For a cell $c$ of size $R_c$ at distance $r_{ic}>R_c/\theta$, the order-$p$ multipole
approximation satisfies
\begin{equation}
\big\|\hat{a}_{c\to i}^{(p)}-a_{c\to i}^{\mathrm{exact}}\big\|
\le C_p\,M_c\,\frac{R_c^{\,p+2}}{r_{ic}^{\,p+4}}.
\end{equation}
For the monopole ($p=0$) this is the familiar $C_0 M_c R_c^{2}/r_{ic}^{4}$: relative to the
monopole force $M_c/r_{ic}^{2}$ it is exactly of order $(R_c/r_{ic})^{2}$. This is the
standard remainder of a multipole expansion about the centre of mass
\citep{greengard1987}, and it is \emph{derived} below (``Derivation of A3'') rather than
asserted; that derivation also yields the explicit factor
$C_p\lesssim(2p+5)/(1-\theta)$. Note that the
expansion is taken about the cell centre of mass, so the \emph{dipole} moment vanishes
identically and the first non-vanishing omitted term is the quadrupole; this is why the
monopole's leading remainder carries $R^2$ and not $R$.

\subsection{Step 1: far-field truncation error at a single node}
Summing A3 over the far set $\mathcal{C}_i$ and factoring out the opening ratio,
\begin{align}
\big\|a_i^{\mathrm{far}}-a_i^{\mathrm{far,exact}}\big\|
&\le\sum_{c\in\mathcal{C}_i}C_p M_c\frac{R_c^{\,p+2}}{r_{ic}^{\,p+4}}\nonumber\\
&=C_p\sum_{c\in\mathcal{C}_i}\frac{M_c}{r_{ic}^{2}}\Big(\frac{R_c}{r_{ic}}\Big)^{p+2}\nonumber\\
&\le C_p\,\theta^{\,p+2}\sum_{c\in\mathcal{C}_i}\frac{M_c}{r_{ic}^{2}},
\label{eq:step1}
\end{align}
where the last inequality uses the opening criterion $R_c/r_{ic}<\theta$. The near field is
exact by construction, so the total single-node operator error is bounded by the far-field
term:
\begin{equation}
\|a_i^{\mathrm{phys}}-a_i^{\mathrm{exact}}\|
\le C_p\,\theta^{\,p+2}\sum_{c\in\mathcal{C}_i}\frac{M_c}{r_{ic}^{2}}.
\label{eq:step1b}
\end{equation}

\subsection{Step 2: from opening angle to $R/r$}
Two lengths appear from here on and they are \emph{not} the same symbol, so we fix them. Write
$r_{ic}$ for the distance from node $i$ to a particular far cell $c$, and let
$r:=\min_{c\in\mathcal C_i} r_{ic}$ be the \emph{closest} such distance. Every far cell
satisfies $R_c/r_{ic}<\theta$, and writing $R$ for a common upper bound on the cell size, the
remaining sum obeys $\sum_{c\in\mathcal{C}_i}M_c/r_{ic}^{2}\le M_{\rm tot}/r^{2}$
\emph{with $r$ the closest distance} --- it would not hold for an average $r_{ic}$. We flag that
this step is deliberately loose: whenever the far cells span a wide range of distances it
replaces the true (smaller) sum by its worst case, which is one reason $C_p$ stays symbolic
rather than tight. Here $M_{\rm tot}$ is the total mass in the far set, and we obtain
\begin{equation}
\|a_i^{\mathrm{phys}}-a_i^{\mathrm{exact}}\|
\le C_p\,\frac{M_{\rm tot}}{r^{2}}\Big(\frac{R}{r}\Big)^{p+2}.
\label{eq:step2}
\end{equation}
Relative to the monopole force scale $M_{\rm tot}/r^{2}$ this is exactly the factor
$(R/r)^{p+2}$, which is the form quoted in Proposition~\ref{prop:bound}. The prefactor $M_{\rm tot}/r^{2}$ is
problem-dependent but independent of the network, so we absorb it into $C_p$ and write
\begin{equation}
\|a_i^{\mathrm{phys}}-a_i^{\mathrm{exact}}\|\le C_p\Big(\frac{R}{r}\Big)^{p+2}\|h\|,
\label{eq:step2b}
\end{equation}
where the feature-scale factor $\|h\|$ is carried along for dimensional consistency with the
layer-wise propagation below.

\subsection{Step 3: propagation through $K$ layers}
Let $e_i^k=h_i^k-h_i^{k,\mathrm{exact}}$ denote the feature error at layer $k$. Each layer
applies
\begin{equation}
h_i^{k+1}=\phi\big(h_i^k,\;\mathrm{AGG}(\{h_j^k\}_{j\in\mathcal{N}(i)}),\;a_i^{\mathrm{phys}}\big).
\end{equation}
Applying A1 and A2 and isolating the operator error
$\epsilon_i:=\|a_i^{\mathrm{phys}}-a_i^{\mathrm{exact}}\|$,
\begin{equation}
\|e_i^{k+1}\|\le\ell\big(\alpha\max_j\|e_j^k\|+\epsilon_i\big).
\label{eq:step3a}
\end{equation}
Taking the max over nodes, $\|e^{k+1}\|\le\beta\|e^k\|+\ell\epsilon$ with $\beta:=\ell\alpha$
and $\epsilon:=\max_i\epsilon_i$. Unrolling from $k=0$ (where $e^0=0$, since inputs are
identical),
\begin{equation}
\|e^K\|\le\ell\epsilon\sum_{k=0}^{K-1}\beta^k
=\ell\epsilon\,\frac{1-\beta^K}{1-\beta},
\label{eq:step3b}
\end{equation}
which converges under A2's effective-contraction condition $\beta=\ell\alpha<1$.

\subsection{Step 4: combining Steps 2 and 3}
Substituting \eqref{eq:step2b} into \eqref{eq:step3b} and absorbing the Lipschitz factor
$\ell$ into the normaliser $\gamma$,
\begin{equation}
\big\|\mathrm{HMMP}(G)-\mathrm{Exact}(G)\big\|
\le\gamma\,C_p\Big(\frac{R}{r}\Big)^{p+2}\|h\|\,\frac{1-\beta^K}{1-\beta}.
\label{eq:theorem}
\end{equation}
This is the statement of Proposition~\ref{prop:bound}. $\square$

\subsection{Derivation of A3 (the multipole remainder)}
\label{app:a3}
A3 is the standard remainder of a multipole expansion taken about the centre of mass. We
derive it explicitly because the index shift is easy to get wrong, and because an earlier
draft simply asserted it.

\emph{Step A (which potential term is omitted).} Shift the origin to the cell's centre of mass,
so $\sum_a m_a y_a=0$ and $|y_a|\le R_c$. For $r=|x|>R_c$ the standard expansion gives
\begin{equation}
\begin{aligned}
\sum_a \frac{m_a}{|x-y_a|}&=\sum_{n\ge 0}\frac{Q_n(\hat x)}{r^{n+1}},\\
Q_n(\hat x)&:=\sum_a m_a\,|y_a|^n P_n(\cos\gamma_a),
\end{aligned}
\end{equation}
with $\gamma_a$ the angle between $x$ and $y_a$. Then $Q_0=M_c$ and
$Q_1=\hat x\cdot\sum_a m_a y_a=0$ \emph{identically}. Truncating at order $p$ therefore retains
$n\le p+1$ (the $n=1$ term contributes nothing), so the first \emph{omitted} term is $n=p+2$
and the potential remainder is
\begin{equation}
\begin{aligned}
\mathcal R_p(x)&=\sum_{n\ge p+2}\frac{Q_n(\hat x)}{r^{n+1}},\\
|Q_n(\hat x)|&\le\sum_a m_a|y_a|^n\le M_c R_c^{\,n},
\end{aligned}
\end{equation}
using $|P_n|\le1$. This is the index shift: because the dipole vanishes, ``order $p$'' first
omits $R^{p+2}$, not $R^{p+1}$ --- which is why the monopole's leading remainder carries $R^2$.

\emph{Step B (from potential to force).} The acceleration is $a=-\nabla\Phi$ and the potential
error is $\Phi_{\mathrm{err}}=-G\,\mathcal R_p$, so the force error is $G\,\nabla\mathcal R_p$.
Differentiating one term $h_n(x)=Q_n(\hat x)/r^{n+1}$,
\begin{equation}
\begin{aligned}
\partial_k h_n&=\frac{\partial_k Q_n(\hat x)}{r^{n+1}}
  -(n+1)\frac{Q_n(\hat x)\,\hat x_k}{r^{n+2}},\\
\big|\partial_k Q_n(\hat x)\big|&\lesssim \frac{n\,M_c R_c^{\,n}}{r},
\end{aligned}
\end{equation}
where the second estimate uses that $Q_n$ is a spherical harmonic of degree $n$ in $\hat x$, so
its angular gradient scales as $n$. Hence
$|\partial_k h_n|\lesssim (2n+1)M_c R_c^{\,n}/r^{n+2}$. Summing over $n\ge p+2$ with
$t:=R_c/r<\theta$ gives a geometric series whose \emph{coefficients also grow}, so it may
\emph{not} be bounded by its first term times $1/(1-t)$: with $N:=p+2$,
\begin{equation}
\sum_{n\ge N}(2n+1)t^{\,n}
=\frac{t^{\,N}}{1-t}\Big[(2N+1)+\frac{2t}{1-t}\Big],
\end{equation}
so dropping the second bracket term would understate the bound. Hence
\begin{equation}
\big\|\nabla\mathcal R_p\big\|\;\le\;
\underbrace{C_p}_{\textstyle \lesssim\,\frac{2p+5}{1-\theta}+\frac{2\theta}{(1-\theta)^2}}\;
M_c\,\frac{R_c^{\,p+2}}{r^{\,p+4}},
\end{equation}
which is A3, with an explicit dependence on both $p$ and $\theta$. The pole is second order,
not first, and at $\theta=0.5$, $p=0$ the bracket raises the constant from $10$ to $14$.

\emph{The constant is sharp, and it is not $(2p+5)$.} The per-order factor can be improved from
$2n+1$ to $n+1$ by bounding the angular and radial parts of $\nabla h_n$ jointly instead of
separately: with $S_n(z)=|y|^n r^{-(n+1)}P_n(\hat z\cdot\hat y)$ one has the exact identity
$(\hat y\cdot\nabla)S_n=-(n+1)|y|^n r^{-(n+2)}P_{n+1}$, so $|\nabla S_n|\le(n+1)|y|^n r^{-(n+2)}$,
attained at $\hat y=\pm\hat z$. Applying the same argument one order up and reducing the resulting
tensor with the Legendre equation gives the second-derivative constant exactly:
$|\nabla^2S_n|_F\le\sqrt{3/2}\,(n+1)(n+2)\,|y|^n r^{-(n+3)}$, also sharp, the $\sqrt{3/2}$
being an exact eigenvalue ratio and not a fit (Appendix~\ref{app:secondorder}). Both constants are checked numerically against finite differences for $l\le12$ (\texttt{scripts/\allowbreak \_\allowbreak hessian\_\allowbreak constant.\allowbreak py}, agreement to $2.8\times10^{-7}$), and the bound built on them is tested where a too-small constant would fail loudly: \textbf{zero violations over $13{,}312$ sampled nodes} --- six float64 runs at $2.3$M and one at $16.7$M, under all three acceptance criteria (Appendix~\ref{app:secondorder}). \emph{An earlier revision stated a symbolic verification to $n=40$ by a script that is not in the repository; we do not assert it.}

\emph{What remains assumed.} The factor $(2p+5)/(1-\theta)$ is explicit, but the $\lesssim$
hides (i)~the angular-gradient constant and (ii)~the moment bound
$\sum_a m_a|y_a|^{p+2}\le M_cR_c^{p+2}$, which is loose for a cell whose mass concentrates near
its centre. We therefore keep $C_p$ symbolic throughout and test only the \emph{scaling} in
$\theta$ and in $p$ --- which is exactly what Section~\ref{sec:order} measures.

\subsection{Remarks}

\paragraph{Remark B.1 (monopole case, and what is measured).}
With $p=0$, $(R/r)^{p+2}=(R/r)^2$, recovering the form quoted in the main text. Retaining
the quadrupole ($p=1$) gives $(R/r)^3$, i.e.\ a faster decay. \textbf{This order dependence
is measured directly} in Section~\ref{sec:order}: the complete operator's error ratio between $p{=}0$ and
$p{=}1$ is proportional to $1/\theta$ with a constant stable to $4.93\%$ over $21$
$(N,\theta)$ settings on controlled synthetic clouds, which is the $(R/r)$ factor that distinguishes
the two orders. On the real snapshot that constant is not stable --- it varies by $1.6$--$1.9\times$
over the same range of $\theta$, and its tidal counterpart is not resolvable at $2.3$M
(Table~\ref{tab:camelsorder}) --- so what is established on clouds is the scaling, and what transfers
to real data is the sign of the shift. What is \emph{not} measured is the layer-count factor or the
absolute constant $C_p$.

\paragraph{Remark B.2 (scope).}
The bound covers the far-field multipole truncation only. The near-leaf sums and the $k$-NN
edge sums are computed exactly and contribute no truncation error. The bound therefore
controls the physical operator $a^{\mathrm{phys}}$, \emph{not} the downstream
density-regression MSE; the latter also depends on the learned message function and the
update MLP.

\paragraph{Remark B.3 (constants are not fitted).}
The constants $\ell,\alpha,C_p,\gamma$ enter as upper-bound factors. We do not fit them to
the measured errors. What the experiments actually test is the bound's functional form: its
dependence on the \emph{opening angle} ($R^2/r^2$ behaviour in $\theta$) and on the
\emph{multipole order} (the $1/\theta$ ratio between $p{=}0$ and $p{=}1$, Section~\ref{sec:order}). The
layer-count factor $(1-\beta^K)/(1-\beta)$ is not isolated by any experiment, and a fully
quantitative bound-versus-measurement ratio is therefore not reported.

\paragraph{Remark B.4 (what $C_p$ now contains).}
In \eqref{eq:step2} the single-node error carries the explicit prefactor
$M_{\rm tot}/r^{2}$, i.e. the monopole force scale. Absorbing it into $C_p$ does not affect
the bound's dependence on $p$, $\theta$, or $K$ --- the whole point of the derivation is
that those enter only through $(R/r)^{p+2}$ and the geometric factor. \textbf{The consequence
for $C_p$:} it now silently contains a mass-distribution- and geometry-dependent factor that
we make no attempt to bound tightly, so $C_p$ is a symbol for "some constant of the mass
distribution and the tree", not a computable number. This is precisely why the experiments
verify the \emph{scaling} (the dependence on $\theta$ and on the order $p$) and not the
\emph{constant}: an absolute bound-versus-measurement comparison would require bounding
$M_{\rm tot}/r^2$ and the cell moments, which we have not done.

\paragraph{Remark B.5 (what is a sketch).}
Steps 1--4 use the multipole remainder bound A3, which \textbf{is derived} above
(Appendix~B, ``Derivation of A3'') rather than asserted; what that derivation leaves loose is
the moment bound $\sum_a m_a|y_a|^{p+2}\le M_cR_c^{p+2}$ and the angular-gradient constant,
which is why $C_p$ stays symbolic. We also do not prove that the \emph{trained} network
satisfies the effective-contractivity hypothesis $\beta<1$; it is a sufficient condition we
support empirically, not a certified property of the learned weights. So the accurate
description of the result is: a bound whose \emph{form} is proved under A1--A3 and whose
scaling in $\theta$ and $p$ is verified, but whose constants are not evaluated.

\section{The second-derivative error law}\label{app:secondorder}
A3 bounds the \emph{force} remainder. The quantity a tidal-field or adjoint application consumes is
the second derivative, and the same expansion can be bounded one order up. This appendix states the
law, gives the exact per-order constants, and reports the measured looseness of the resulting bound
on controlled single cells. Everything here is per \emph{interaction}; Section~\ref{sec:order} is the
aggregate measurement, and the two must not be read as the same object.

\noindent\textbf{How this differs from the classical potential remainder, term by term.} The classical
FMM error analysis~\citep{cheng1999fast} states four remainders, all for the \emph{potential} $\Phi$:
Theorem~2.1, eq.~(7),
$|\Phi-\sum_{n\le p}\cdots|\le\bigl(\sum_i|q_i|/(r-a)\bigr)(a/r)^{p+1}$, for sources inside a sphere
of radius $a$ about the expansion centre; Theorem~2.4, eq.~(18),
$\le\bigl(\sum_i|q_i|/(ca-a)\bigr)(1/c)^{p+1}$ for the multipole-to-local conversion; Theorem~2.2,
eq.~(10) for the local expansion; and Observation~3.4, eq.~(41), $|\Phi-\Phi_\varepsilon|<A\varepsilon$
for the exponential (plane-wave) representation. Theorem~2.3, eq.~(15) supplies the multipole
translation. \textbf{No force remainder and no derivative remainder appears in that paper} ---
``derivative'', ``gradient'' and ``Hessian'' do not occur in its 31 pages
(Appendix~\ref{app:provenance}, D7).

Three consequences, each of which matters for how the claim is read.

\emph{(i) The missing object is the twice-differentiated one, but it is not their bound rescaled.}
Raising a multipole remainder by one derivative raises the power of $t$ and of $1/r$ by one, so the
object sought has the shape $t^{l+1}$ with one further power of $1/r$ per derivative. It does not
follow by reusing their inequality, for the reason given in Section~\ref{sec:limits}: with the
expansion about the centre of mass the dipole vanishes identically, so the lowest omitted order is
fixed rather than free, and the naive one-order-per-derivative extension of a \emph{force} remainder
diverges as $\theta\to0$ at $p{=}0$. The derivation has to be redone, which is what
Section~\ref{sec:order} and this appendix do.

\emph{(ii) Their $a$ is the same kind of length as ours; what differs is what it is compared
against.} In Theorem~2.1, $a$ is the radius of a sphere about the expansion centre containing the
sources. When that centre is the centre of mass, such a sphere has radius exactly $\max_b|y_b|$ ---
the $a$ of \eqref{eq:cellbound}. So the geometry distinction claimed in the main text is \emph{not} a
difference in the definition of this length: it is that the tree's acceptance criterion is set by the
cell's geometric size $R$, while the bound is set by $a$, and $a/R$ is a property of the mass
distribution that reaches $\sqrt3$ for a cell whose mass sits in one corner
(Section~\ref{sec:criteria}). Their theorem does not connect $a$ to any acceptance criterion, so it
cannot see that gap; \eqref{eq:cellbound} with the dilation of Section~\ref{sec:criteria} does.

\emph{(iii) The bases differ, so the constants are not directly comparable.} Their expansion is in
spherical harmonics with plane-wave (diagonal) translation operators; the bound here is in Legendre
harmonics with direct differentiation, and $\mu_l$ and $\Lambda_l$ of
\eqref{eq:lem1}--\eqref{eq:lem2} are the exact operator norms \emph{for that basis}. This narrows
rather than broadens what is claimed: even inside the classical FMM error-analysis framework, no
order-explicit constant for the second-derivative remainder has been written down.

\subsection{The two lemmas}

Write the $l$-th angular moment of a cell about its centre of mass as
\begin{equation}
F_l(\hat x):=\sum_b m_b\,|y_b|^l P_l(\hat x\cdot\hat y_b),
\qquad
h_l(x):=r^{-(l+1)}F_l(\hat x),
\end{equation}
so that the exact potential is $\sum_l h_l$ and truncating at order $L$ leaves
$\mathcal R_L=\sum_{l>L}h_l$. Then, with $\|\cdot\|$ the $L^2$ norm on the sphere normalised by
$1/4\pi$ and $|\cdot|_F$ the Frobenius norm,
\begin{align}
\big\|\nabla h_l\big\| &\le \mu_l\,\|F_l\|, &
\mu_l &= (2l+1)\sqrt{l+1},
\label{eq:lem1}\\
\big\|\nabla^2 h_l\big\|_F &\le \Lambda_l\,\|F_l\|, &
\Lambda_l &= (2l+1)\sqrt{(l+1)(l+2)(2l+3)} .
\label{eq:lem2}
\end{align}

\emph{Derivation.} Expand $F_l=\sum_m c_m\tilde Y_{lm}$ in a basis orthonormal under
$d\Omega/4\pi$, so $\|F_l\|^2=\sum_m|c_m|^2$. Since $\nabla^2h_l=\sum_m c_m\nabla^2 h_{lm}$, Cauchy--Schwarz
gives $|\nabla^2h_l(\hat x)|_F\le\|F_l\|\sqrt{\sum_m|\nabla^2h_{lm}(\hat x)|^2_F}$, and the sum is
\emph{independent of $\hat x$} by rotation invariance, so it is one number per order. Evaluating it
uses the reproducing kernel of the degree-$l$ space, which by the addition theorem is
$\sum_m\tilde Y_{lm}(\hat x)\tilde Y_{lm}(\hat x')=(2l+1)P_l(\hat x\cdot\hat x')$:
\begin{equation}
\Lambda_l^2=\sum_{ij}\partial_i\partial_j\partial'_{i}\partial'_{j}
\Big[(2l+1)P_l(\hat x\cdot\hat x')\,|x|^{-(l+1)}|x'|^{-(l+1)}\Big]_{x=x'} .
\end{equation}
At $x=x'$ the only available tensors are $\delta_{ij}$ and $\hat x_i\hat x_j$; contracting and
reducing with the Legendre equation gives $\Lambda_l^2=(2l+1)^2(l+1)(l+2)(2l+3)$ and, one order down,
$\mu_l^2=(2l+1)^2(l+1)$. The $l{=}0$ case is exact and checkable by hand: $\Lambda_0=\sqrt6$ and
$\|\nabla^2 r^{-1}\|_F=\sqrt6/r^3$ for $F_0\equiv1$. Note $\Lambda_0\ne0$: an exterior solid harmonic
is harmonic, so only its \emph{trace} vanishes, not its Hessian. Both constants were verified
independently by fourth-order finite differences of $h_{lm}$ over the full real basis for $l\le12$,
agreeing with the closed forms to $2.8\times10^{-7}$ (\texttt{scripts/\allowbreak \_\allowbreak hessian\_\allowbreak constant.\allowbreak py}).

\subsection{The cell bound, and its geometry dependence}

Summing \eqref{eq:lem2} over the omitted orders and using $\|F_l\|\le\sup|F_l|\le\sum_b m_b|y_b|^l=:S_l$,
\begin{equation}
\big\|\nabla^2\mathcal R_L\big\|_F
\;\le\;\sum_{l>L}\Lambda_l\,S_l\,r^{-(l+3)}
\;\le\;\frac{M}{r^{3}}\sum_{l>L}\Lambda_l\,t^{\,l},
\qquad t:=\frac{a}{r},\quad a:=\max_b|y_b| .
\label{eq:cellbound}
\end{equation}
This is Dehnen-shaped but one derivative higher: the order factor $t^{L+1}$ is the same one Dehnen's
force bound carries, there is one extra power of $1/r$, and because $\Lambda_l\simeq2\sqrt2\,l^{5/2}$
grows faster than Dehnen's linear-in-$l$ coefficient the pole in $1/(1-t)$ is of higher order.
Bounding $\Lambda_l\le2(2l+1)(l+1)(l+2)$ makes the sum elementary and gives the explicit form
$\|\nabla^2\mathcal R_L\|_F\lesssim2(M/r^3)t^{L+1}(L+1)^3/(1-t)^4$ in the large-$L$, $t\to1$ regime.

\noindent\textbf{The criterion's $\theta$ is not the bound's $t$, and the gap is the cell's geometry.}
The acceptance test bounds $R/r$ with $R$ the cell's geometric size, while \eqref{eq:cellbound} needs
$t=a/r$ with $a$ the distance from the \emph{centre of mass} to the farthest particle. The two differ
by the dilation $(a/R)^{l_{\min}}$, where $l_{\min}$ is the first non-vanishing omitted order --- $2$
for the monopole, because the dipole dies at the centre of mass, and $3$ once the quadrupole is
retained. $a/R$ is a property of the mass distribution, not of $\theta$: measured here it is $0.837$
for a uniform cube, $0.698$ for a slab of aspect $1/4$ and $0.693$ for one of aspect $1/8$, while a
cell whose mass sits in one corner reaches $\sqrt3=1.732$. So a size-based opening test at fixed
$\theta$ is \emph{not} a fixed error bound: the same $\theta$ admits errors differing by up to $3\times$
at $l_{\min}=2$ and $5.2\times$ at $l_{\min}=3$ purely from where the mass sits inside the box. This
is the quantitative form of the geometry dependence that defect D8
(Appendix~\ref{app:provenance}) exposed on the octree side.

\subsection{Measured looseness, and one recorded failure}

Table~\ref{tab:secondorder} reports \eqref{eq:cellbound} against the exact single-cell remainder: a
cell of 1500 unit masses in a box of side $R=1$, queries on 32 directions per radius, three seeds,
and the reference a brute-force sum over the same particles. \textbf{The bound is never violated.}
Its looseness is $11.5$--$48.5\times$ depending on radius and order, and it is stable in $r$, which is
the diagnostic that the \emph{power} is right and only the constant is conservative --- the opposite
signature from a bound whose exponent is wrong, which would drift systematically in $r$.

\begin{table}[t]
\centering\small
\caption{Bound~\eqref{eq:cellbound} over the measured single-cell second-derivative remainder.
``harmonic'' feeds the exact $\|F_l\|$; ``isotropic'' uses the crude $S_l\le Ma^l$, i.e.\ it throws
away the cell's symmetry. Both are valid; the ratio is how much is thrown away. \textbf{No sampled
interaction violates the bound.} The last column is the first version of the constant,
$\sqrt{3/2}(l+1)(l+2)r^{-(l+3)}$, which is the maximum over the \emph{zonal} harmonic rather than the
operator norm on general harmonics: it sits \emph{below} the measured remainder, so it was never a
bound. It is kept on the record rather than deleted. All entries measured, none fitted.}
\label{tab:secondorder}
\resizebox{\ifdim\width>\linewidth\linewidth\else\width\fi}{!}{%
\begin{tabular}{llrrrr}
\toprule
cell & retained & $r/R$ & bound (harmonic) / meas. & bound (isotropic) / meas. & earlier (non-)bound / meas. \\
\midrule
cube & $l\le0$ & 1.5  & 48.54 & 1428.54 & 10.67 \\
cube & $l\le0$ & 4.0  & 19.92 & 1203.23 & 5.56 \\
cube & $l\le0$ & 32.0 & 11.54 & 5479.01 & 3.68 \\
cube & $l\le2$ & 1.5  & 48.08 & 1261.83 & 9.96 \\
cube & $l\le2$ & 32.0 & 27.50 & 1647.20 & 6.86 \\
slab $1/8$ & $l\le0$ & 1.5  & 19.34 & 144.80 & 5.30 \\
slab $1/8$ & $l\le0$ & 32.0 & 11.62 & 527.83 & 3.72 \\
slab $1/8$ & $l\le2$ & 1.5  & 41.71 & 476.56 & 8.77 \\
slab $1/8$ & $l\le2$ & 32.0 & 31.78 & 1762.80 & 7.72 \\
\bottomrule
\end{tabular}
}
\end{table}

\noindent\textbf{Why the isotropic form is unusable, and what that says about symmetric cells.} The
$S_l\le Ma^l$ column is loose by up to $5479\times$, and the looseness \emph{grows with $r$} --- an
apparent drift of $r^{+0.6}$ that is pure book-keeping. The cause is symmetry: a uniform cube has
inversion symmetry about its centre of mass, so its odd moments vanish identically while
$S_3,S_5,\dots$ stay finite, and the sum then carries orders that contribute nothing to the true
remainder. Feeding the exact $\|F_l\|$ instead makes those orders vanish as they should and collapses
the looseness by two orders of magnitude. A practitioner should therefore never use the
extent-only form; the harmonic norm costs $O(n(2l+1))$ per order and no pairwise work.

\noindent\textbf{What this bound is and is not.} It is a \emph{provable} per-interaction upper bound
on the second-derivative remainder of a truncated multipole expansion about a cell's centre of mass,
with the order $L$, the size-to-distance ratio $t$ and the cell's mass distribution and aspect ratio
entering explicitly, and with constants that are exact rather than fitted. It is not tight: an
$11$--$48\times$ factor separates it from the measured remainder on these cells, so it certifies
rather than predicts. \textbf{On real data it is looser still, and the looseness is not a constant.}
Running the same test on the actual admitted cells of the CAMELS $L25n256$ snapshot (a $2\times10^{4}$
subsample, $128$ query nodes, brute-force reference, $\theta=0.5$, quadrupole retained) gives
\textbf{no violation on any node}, with a looseness of \textbf{median $306\times$}, tightest $11\times$
and worst $1026\times$; the corresponding figures are $183\times$ on a synthetic cloud at
$N=4000$ and $267\times$ at $N=50000$. \textbf{At the full $2.3$M scale}, over the $1{,}643{,}251$ far
interactions actually admitted at $\theta=0.5$ with the centre-of-mass criterion and $2{,}048$ sampled
query nodes, the bound again holds on \textbf{every node --- $0$ violations} --- with a looseness of
\textbf{median $426\times$} (tightest $30\times$, worst $2534\times$).

\textbf{That is a sampled statement, so we hardened it rather than trusting one draw.} The per-node
remainder is heavy-tailed --- established independently in Section~\ref{sec:scale} --- so a bound
verified on one draw of $2{,}048$ nodes is a claim about that draw until it is repeated. We repeated
the $2.3$M test in float64 under \emph{all three} acceptance criteria and \emph{five} independent
query sets: six runs, $12{,}288$ sampled nodes in total, and \textbf{$0$ violations in every one},
with the tightest ratio anywhere $20.5$ and the median looseness $407$--$513\times$. The geometry
premises are re-audited per run and hold throughout ($a/r\le0.750$, $a/\text{size}\le1.245$ against
the invariant's $\sqrt3$, and no interaction with $a\ge r$). \textbf{At the paper's title scale} we also ran
the same test on the entire $16.7$M snapshot --- deliberately the same $N$ as the operator validation
of Section~\ref{sec:scale}, so that bound and operator are verified at the same scale rather than at
two different ones: $3{,}927{,}827$ cells, $1{,}060{,}190$ far interactions actually admitted, $1{,}024$
sampled query nodes, and again \textbf{$0$ violations}, with a looseness of
\textbf{median $436\times$} (tightest $28.3\times$, worst $3671\times$) and the same geometry premises
holding ($a/r\le0.578$, $a/\text{size}\le1.274$, no interaction with $a\ge r$, no non-convergent
cell). \textbf{So the bound is verified at the scale in this paper's title, not only at the scale of
its tables.} The looseness chain is $186\times\to306\times\to426\times\to436\times$ across
$4$k$\to20$k$\to2.3$M$\to16.7$M: it grows with scale and appears to be \emph{saturating}, which is the
signature of a bound whose exponent is right --- a bound with the wrong power of $r$ would drift
without limit instead.

\noindent\textbf{The measured side of that test must be evaluated in float64, and getting this wrong
manufactures violations.} The remainder is a \emph{difference} of two gradient sums whose norms are
$\sim10^{4}$ times the difference itself, so a float32 evaluation is dominated by rounding: at the
median it inflates the measured remainder by $13.9\times$, which simultaneously \emph{understates}
the looseness and invents violations. Concretely, the identical $2.3$M configuration in float32
reports \textbf{two violations} and a looseness of $32.6\times$; in float64 the same run reports
\textbf{none} and $425.6\times$ --- with every bound-side quantity agreeing to six significant
figures and the tail share bit-identical. \texttt{scripts/\allowbreak error\_\allowbreak law\_\allowbreak bound\_\allowbreak sampled.\allowbreak py
--field-dtype float64} is the committed configuration; the float32 numbers are kept in
Appendix~\ref{app:provenance} as the reason the flag exists. That is
why the same machinery is paired in
Section~\ref{sec:selfcertify} with an a posteriori estimator built from the operator's own
next-order surplus, which is far tighter because it is the omitted term itself rather than a bound on
it --- and why that estimator's failure, reported there, matters. We are also unable to claim that no
derivative remainder exists in the classical FMM literature
(Appendix~\ref{app:provenance}, D7); what is new here is that the constant is exact, that the geometry
dependence is explicit, and that the bound is verified against measured per-interaction remainders
rather than asserted.

\section{Reading the cost tables}\label{app:scaling}
\noindent\textbf{The opening-angle sweep, and what it is not.} A separate sweep varies the
\emph{opening angle} $\theta$ at fixed multipole truncation (monopole only, $p=0$), over the six
values $\theta=0.467,0.35,0.28,0.233,0.20,0.175$; the measured relative force error falls
monotonically as the criterion tightens: $0.0032450$, $0.0014542$, $0.00077348$, $0.00047959$,
$0.00028759$, $0.00018170$. (The pre-D8 artifact stored $0.0312507$ for the loosest setting, i.e.\
$9.6\times$ worse: this curve is another measurement that moved with the geometry fix.) The
comparison curve stored alongside is a \emph{scaling proxy}, not Proposition~\ref{prop:bound}
evaluated with realised constants, so what this supports is narrow: the error decreases with a
tighter opening, consistent with the $R/r$ dependence. It does \textbf{not} measure the \emph{order}
dependence, because $p$ never changes here --- that is Section~\ref{sec:order}, and the two sweeps
must not be conflated. The artifact stores this $\theta$-sweep under the key \texttt{error\_vs\_order}
although $\theta$ is the swept variable; the misnomer is flagged in Appendix~\ref{app:provenance} so a
reader of the artifacts is not misled.

The tables in this appendix and in Section~\ref{sec:scaling} measure three different quantities, and
they must not be combined into a single curve:
\begin{itemize}\itemsep2pt
\item \textbf{Forward-pass time} (Table~\ref{tab:time}) is a single forward pass on a controlled
particle set of increasing $N$ at a fixed configuration. It is the only time measurement here that
varies $N$ cleanly.
\item \textbf{Peak memory} (Tables~\ref{tab:mem} and~\ref{tab:memdeep}) is the peak CUDA allocation
of a single forward pass. The capped and uncapped interaction-list series are \emph{different
configurations}, so the two are not one function of $N$.
\item \textbf{Opening-angle error} (Section~\ref{sec:scaling}) is a force-accuracy measurement, not a
cost measurement.
\end{itemize}
\noindent None of these is a full-simulation wall-clock. We deliberately report no such number: the
operator's cost inside a production loop depends on how often the particle set changes, and therefore
how often the cached tree and interaction list must be rebuilt --- which we have not benchmarked. The
caching result is a statement about repeated calls on a \emph{static} particle set.

\section{Forward-pass time table}\label{app:tables}
\begin{table}[h]
\centering
\caption{Forward-pass time of the HMMP operator versus $N$ at fixed configuration
(\texttt{scripts/\allowbreak scaling.\allowbreak py}), post-D8, one A800. Over a $20\times$ range in $N$ the time grows
$47.79\times$. We do \textbf{not} read an exponent from this: the same configuration on CPU grows
$26.68\times$ over the same range, so the ratio is device-dependent and the small-$N$ end is
overhead-dominated. The admitted far-interaction count --- measured by the two-arm comparison of
\texttt{scripts/\allowbreak compare\_\allowbreak d8\_\allowbreak cost\_\allowbreak ab.\allowbreak py}, which prints the two reference ratios beside it --- grows
$71.0\times$ over the same range, i.e.\ faster than $O(N\log N)$, because the benchmark clouds are
not self-similar. Four points cannot separate these, so we report the measurements rather than a
fitted exponent.}
\label{tab:time}
\resizebox{\ifdim\width>\linewidth\linewidth\else\width\fi}{!}{%
\begin{tabular}{lrr}
\toprule
$N$ & HMMP time (s) & ratio vs $N{=}1000$ \\
\midrule
1{,}000 & 0.015203 & 1.000000 \\
3{,}000 & 0.050825 & 3.343671 \\
10{,}000 & 0.223754 & 14.720343 \\
20{,}000 & 0.726464 & 47.792956 \\
\bottomrule
\end{tabular}
}
\end{table}

\noindent The pre-D8 table gave $0.205239\to4.155135$\,s, a $20.25\times$ growth: the geometry fix
made the operator admit more far interactions and therefore cost more, so the two tables are not the
same measurement and the earlier one must not be quoted beside this. The brute-force reference in the
A800 run is \textbf{not quotable} --- it is non-monotone ($3.4$\,ms at $N=1{,}000$ against $1.1$\,ms at
$N=3{,}000$), i.e.\ overhead-dominated at both points --- so we quote the CPU harness instead:
$0.039158$\,s at $N=1{,}000$ and $0.265796$\,s at $N=3{,}000$, against the tree's CPU times of
$0.260325$\,s and $0.886076$\,s at the same $N$. The tree is still the slower method at both points,
which is why no crossover between the tree and the direct sum is quoted: we can see it has no
advantage yet at $N=2{,}000$ (Section~\ref{sec:exact}) and that it scales better, but we have not
measured where the two meet.

\section{Memory-scaling tables}\label{app:memtable}
\begin{table}[h]
\centering
\caption{Peak CUDA allocation for a single forward pass, hidden $64$, interaction list
\emph{uncapped}, post-D8, one A800 (\texttt{scripts/\allowbreak benchmark\_\allowbreak memory.\allowbreak py}). The $2.3$M row failed
under the default caching allocator and completed at $71.42$\,GB under
\texttt{expandable\_segments:True}; it is the only row measured with that setting. Pre-D8 the same
series ended at $11794.97$\,MB.}\label{tab:mem}
\resizebox{\ifdim\width>\linewidth\linewidth\else\width\fi}{!}{%
\begin{tabular}{lr}
\toprule
$N$ & Peak memory (MB) \\
\midrule
10{,}000 & 199.12 \\
50{,}000 & 1014.17 \\
100{,}000 & 2207.70 \\
200{,}000 & 5180.30 \\
500{,}000 & 13064.96 \\
1{,}000{,}000 & 28082.88 \\
2{,}300{,}000 & \textbf{71416.82}$^{\dagger}$ \\
\bottomrule
\end{tabular}
}
\end{table}

\begin{table}[h]
\centering
\caption{Same measurement with the interaction list \emph{capped} at $2.5$M entries, post-D8. This is
a different configuration from Table~\ref{tab:mem}, not a continuation of it; the cap is what the
single-device $2.3$M runs used. Pre-D8 the same series ended at $3780.78$\,MB.}\label{tab:memdeep}
\resizebox{\ifdim\width>\linewidth\linewidth\else\width\fi}{!}{%
\begin{tabular}{lr}
\toprule
$N$ & Peak memory (MB), capped \\
\midrule
10{,}000 & 198.17 \\
50{,}000 & 723.91 \\
100{,}000 & 694.13 \\
200{,}000 & 996.67 \\
500{,}000 & 2481.78 \\
1{,}000{,}000 & 4947.84 \\
2{,}300{,}000 & 11351.59 \\
\bottomrule
\end{tabular}
}
\end{table}

\noindent The capped series is not monotone in $N$ between $50$k and $100$k ($723.91$ against
$694.13$); the difference is $4\%$ in two independent processes and we read it as allocator noise
rather than as a real inversion, but it is small enough to record rather than smooth. $^{\dagger}$The
uncapped $2.3$M row requires \texttt{PYTORCH\_CUDA\_ALLOC\_CONF=expandable\_segments:True}: under the
default allocator the run dies with CUDA out of memory after $72.89$\,GiB is in use
(\texttt{outputs/\allowbreak memory\_\allowbreak hmmp\_\allowbreak 2M3\_\allowbreak uncapped\_\allowbreak postd8.\allowbreak log}, and the passing
\texttt{expandable\_segments} run is \texttt{outputs/\allowbreak memory\_\allowbreak hmmp\_\allowbreak 2M3\_\allowbreak uncapped\_\allowbreak expseg\_\allowbreak postd8.\allowbreak json}).
So the $2.3$M claim survives
the geometry fix, but with much less margin than before: the uncapped configuration is inside one
$80$\,GB device only just, while the capped one is comfortable at $11.35$\,GB. That is a feasibility
statement, not an accuracy statement, and Section~\ref{sec:scaling} says so explicitly.

\section{Aggregation level of the order-shift measurement}\label{app:percell}
\noindent\textbf{The per-$\theta$ means behind Section~\ref{sec:order}'s two-row comparison.}
Synthetic lognormal clouds, post-D8, $p{=}0$ over $p{=}1$, multiplied by $\theta$; per-$\theta$ means
over $N\in\{2{,}000,8{,}000,32{,}000\}$:

\begin{center}
\resizebox{\ifdim\width>\linewidth\linewidth\else\width\fi}{!}{%
\begin{tabular}{lrrrrrr}
\toprule
$\theta$ & 0.15 & 0.20 & 0.25 & 0.30 & 0.40 & 0.50 \\
\midrule
force & 2.532 & 2.474 & 2.402 & 2.352 & 2.336 & 2.249 \\
tidal & 1.836 & 1.839 & 1.759 & 1.743 & 1.740 & 1.726 \\
\bottomrule
\end{tabular}
}
\end{center}

\noindent Over $\theta\le0.5$: force $2.391$, scatter $4.3\%$; tidal $1.774$, scatter $2.9\%$.
Artifact \texttt{outputs/\allowbreak multipole\_order\_sweep\_\allowbreak tidal\_fixed.\allowbreak json}.

Section~\ref{sec:order} reports that the measured order-shift product depends on the particle
distribution, and this appendix gives the measurement behind the qualification that a bound is a
per-interaction object while our numbers are aggregates. On one shared interaction list (real
CAMELS $N=100{,}000$ and a synthetic $N=100{,}000$ cloud, $Q=512$ query nodes, $\texttt{max\_leaf}=8$,
$\theta\in[0.15,0.5]$) we evaluate the $p{=}0$ and $p{=}1$ remainders at three levels, each a ratio
of $p{=}0$ to $p{=}1$ multiplied by $\theta$:

\begin{center}
\resizebox{\ifdim\width>\linewidth\linewidth\else\width\fi}{!}{%
\begin{tabular}{lrrr}
\toprule
level & synthetic & CAMELS & CAMELS drift \\
\midrule
summed per-interaction remainders & $1.169 \pm 0.7\%$ & $1.066 \pm 12.4\%$ & $1.21 \to 0.93$ ($\times1.30$) \\
coherent aggregate over cells & $1.193 \pm 4.0\%$ & $1.037 \pm 22.4\%$ & $1.34 \to 0.81$ ($\times1.65$) \\
per-node relative mean & $1.275 \pm 3.7\%$ & $1.040 \pm 22.7\%$ & $1.36 \to 0.80$ ($\times1.69$) \\
\bottomrule
\end{tabular}
}
\end{center}

\noindent Each cell contribution is the direct sum over \emph{all descendant particles} of the
admitted cell, so both sides of each comparison use the same source set; obtaining that expansion
correctly is what retired an earlier probe that used a leaf-only expansion on cells at every level.
On the controlled cloud the three levels agree and are constant. On the real snapshot all three
drift, the per-interaction level least ($\times1.30$) and the aggregate most ($\times1.65$), so the
drift is not purely an artefact of summing vectors coherently, but it is also not a measurement of a
per-interaction constant: each of the three is a weighted average over the admitted-cell
distribution. Reproduce with \texttt{scripts/\allowbreak diagnose\_\allowbreak aggregate\_\allowbreak vs\_\allowbreak percell.\allowbreak py}
(\texttt{outputs/\allowbreak agg\_\allowbreak vs\_\allowbreak percell\_\allowbreak *\.\allowbreak json}).

\begin{table}[h]
\centering
\caption{The order sweep of Section~\ref{sec:order} repeated on the real CAMELS snapshot at three
scales, post-D8, $\theta\in[0.15,0.60]$, $Q=2{,}048$ query nodes, $\texttt{max\_leaf}=8$. The product
is the consecutive-order error ratio multiplied by $\theta$, i.e.\ the quantity
Proposition~\ref{prop:bound} predicts to be constant; \emph{scatter is the sample standard deviation
over the seven $\theta$ values, relative to the mean}, and ``range'' is the ratio of the largest
product to the smallest. The synthetic value is $2.4$ for the force and $1.8$ for the tidal tensor.
\textbf{The force product is never constant here, and the tidal one stops being measurable at
$2.3$M}: its scatter grows to $53.9\%$ and its product moves over an $8.1\times$ range. What holds at
every scale is the \emph{sign} --- the $p{=}0$/$p{=}1$ ratio exceeds one everywhere ($4.0$--$24.5$
force, $1.4$--$12.1$ tidal).}\label{tab:camelsorder}
\resizebox{\ifdim\width>\linewidth\linewidth\else\width\fi}{!}{%
\begin{tabular}{lrrrrr}
\toprule
$N$ & force product & force scatter & force range & tidal product & tidal scatter \\
\midrule
$10^{5}$ & $3.249$ & $15.7\%$ & $1.64\times$ & $1.758$ & $3.9\%$ \\
$10^{6}$ & $3.743$ & $24.4\%$ & $1.91\times$ & $1.351$ & $35.7\%$ \\
$2.3\times10^{6}$ & $3.318$ & $23.5\%$ & $1.86\times$ & $1.121$ & $53.9\%$ \\
\midrule
\multicolumn{6}{l}{\footnotesize tidal range: $1.10\times$ / $4.09\times$ / $8.10\times$ at the three scales} \\
\bottomrule
\end{tabular}
}
\end{table}

\begin{sloppypar}
\noindent Script \texttt{scripts/\allowbreak sweep\_multipole\_order\_tidal\_camels.\allowbreak py}, artifacts
\texttt{order\_sweep\_tidal\_camels\_n\allowbreak\{100000,1000000,2300000\}\_\allowbreak postd8.json}, whose own summary fields
(\texttt{ratio\_times\_theta\_mean\_*}, \texttt{...\_rel\_scatter\_*}) are the source of the numbers
above. Two remarks a reader should carry away. First, the tidal scatter \emph{grows} with $N$
($3.9\%\to35.7\%\to53.9\%$) while the force scatter does not ($15.7\%\to24.4\%\to23.5\%$), so the
ordering between the two derivatives that holds on controlled clouds and at $10^{5}$ inverts at the
scale this paper is about. Second, at $2.3$M the tidal product at the tightest opening
($\theta=0.15$) is $0.208$, i.e.\ the quadrupole is worth only a factor $1.4$ there against $12.1$ at
the same $\theta$ for $N=10^{5}$: the tidal order shift is not merely noisier at scale, it is weaker.
\end{sloppypar}

\noindent\textbf{The single-cell constant, and whether it survives a holdout.} Section~\ref{sec:order}
asks the aggregate version of this question. The single-cell version is sharper and answerable: take
one source cell, one query point outside it, and the one-order inequality
$\varepsilon_1\le C\rho\varepsilon_0$ with $K=(\varepsilon_0/\varepsilon_1)\rho$ and $C=1/\min K$.
\emph{The inequality as previously reported was not a test}: $C$ was the minimum over the same $900$
samples the inequality was then evaluated on, so it could not be violated there --- it is the tangent
line at the worst sample, and it reproduces exactly ($C=5.520$ tidal, $8.897$ force, over $900$
samples, binding geometry \texttt{G7\_shell}).

So we calibrated on one set and evaluated on another
(\texttt{scripts/\allowbreak e7\_holdout\_bound.\allowbreak py}, \texttt{outputs/\allowbreak e7\_holdout\_bound.\allowbreak json}). Two splits, and
they disagree with each other, which is the point.

\begin{table}[h]
\centering
\small
\caption{Holdout test of the one-order constant. Fold A leaves each cell geometry out in turn:
$C$ is calibrated on the other nine geometries and the inequality is then evaluated on the held-out
one. Fold B calibrates on two random seeds of the three stochastic families and evaluates on a third,
unseen seed. ``inflation'' is the factor by which $C$ must be raised to cover the held-out data, so
$<1$ means the calibration was already conservative there.}\label{tab:e7holdout}
\resizebox{\ifdim\width>\linewidth\linewidth\else\width\fi}{!}{%
\begin{tabular}{lrrrr}
\toprule
held-out set & tidal infl. & tidal viol. & force infl. & force viol. \\
\midrule
\texttt{G2\_two\_mass} & $0.02\times$ & $0/90$ & $0.01\times$ & $0/90$ \\
\texttt{G3\_line} & $0.03\times$ & $0/90$ & $0.01\times$ & $0/90$ \\
\texttt{G4\_triangle} & $0.26\times$ & $0/90$ & $0.13\times$ & $0/90$ \\
\texttt{G5\_uniform} & $0.43\times$ & $0/90$ & $0.23\times$ & $0/90$ \\
\texttt{G6\_bimodal} & $0.09\times$ & $0/90$ & $0.04\times$ & $0/90$ \\
\texttt{G8\_heavy\_centre} & $0.38\times$ & $0/90$ & $0.21\times$ & $0/90$ \\
\texttt{G9\_hollow\_shell\_pair} & $0.38\times$ & $0/90$ & $0.18\times$ & $0/90$ \\
\texttt{G10\_symmetrized\_uniform} & $0.18\times$ & $0/90$ & $0.09\times$ & $0/90$ \\
\midrule
\textbf{\texttt{G7\_shell}} & $\mathbf{2.34\times}$ & $\mathbf{13/90}$ & $\mathbf{4.38\times}$ & $\mathbf{13/90}$ \\
\midrule
fresh seed, stochastic families only & $0.43\times$ & $0/180$ & $0.23\times$ & $0/180$ \\
\bottomrule
\end{tabular}
}
\end{table}

\noindent Three readings. \textbf{(i) The constant is carried by one geometry.} Nine of ten held-out
geometries are covered with room to spare, and the tenth --- the spherical shell, where the octupole
is weakest --- violates the bound at $13$ of $90$ samples and forces $C$ up by $2.3\times$ (tidal) and
$4.4\times$ (force). A constant calibrated on arbitrary clouds therefore does \emph{not} transfer to
the near-degenerate ones. \textbf{(ii) It does transfer between draws.} A fresh seed of the same
families costs nothing ($0$ violations in $180$ unseen samples, and the held-out draws are milder than
the calibration set): the failure is between geometries, not between samples. \textbf{(iii) So the
honest statement of the bound is per family, not universal.} $\sup K$ is finite --- E7's original
question, answered yes --- but a single number $C$ is only valid for a population of cell geometries
that includes its binding member, and the paper reports it that way. This is the single-cell
counterpart of the aggregate drift in Section~\ref{sec:order}, and it explains it: a real snapshot's
admitted cells are a mixture, and its minimum $K$ is set by whichever near-spherical cells it happens
to contain.

\section{Disclosed defects and provenance}\label{app:provenance}
We list here the defects that bear on what this paper reports. Several earlier bookkeeping defects
--- misnamed artifact keys, a table transcribed from an older run, an inert seed control in training
scripts that no result here uses --- are not
about this paper's claims, and are recorded instead in the repository
(\texttt{docs/\allowbreak PROVENANCE.\allowbreak md}), which carries the full evidence chain for every number.

\noindent\textbf{(D2) The quadrupole far-field term was wrong, and is fixed.} It was implemented as
$0.5\,(\hat r^{\top}D\hat r)\,\hat r/r^{3}$, which has the wrong radial power, omits the $D\hat r$
term, drops the trace, and carries the wrong sign; it made the far field roughly $20\times$
\emph{worse} than the monopole alone. The correct form now used is
$0.5\,[-6(D\hat r)-3\,\mathrm{tr}(D)\hat r+15(\hat r^{\top}D\hat r)\hat r]/r^{4}$ with
$D=\sum_j m_j(y_j-c)(y_j-c)^{\top}$. It is pinned by an exact increment identity for two equal
masses, $\Delta a=-6md^{2}/R^{4}$, which the old code missed by a factor of $\sim4000$, and by nine
unit tests. The corrected term is what produces the $2.28$--$2.43\times$ improvement reported in
Section~\ref{sec:order}; without it the quadrupole made results worse, which is the failure mode this
test now prevents.

\noindent\textbf{(D6) The second-derivative reference was not the second derivative, and no quadrupole
contribution to it existed.} \textbf{This invalidated every tidal number this paper reported before
this revision, and is why those numbers changed.} The tidal tensor was computed as
$3d_id_j/r^{2}-\delta_{ij}/r^{3}$, but the Jacobian of $a_i=m\,d_i/r^{3}$ is
$3d_id_j/r^{5}-\delta_{ij}/r^{3}$: the first term was missing a factor $1/r^{3}$. The two expressions
agree only at $r=1$, which is why it survived several revisions. Against
\texttt{torch.autograd} the old expression is wrong by $24{,}000$ at $R=0.05$ (sign included) and by
$2.98$ at $R=5$, while the corrected one matches to $10^{-12}$; on a $2{,}000$-particle cloud the two
tensors differ by $1280\times$ in norm, and the old object was $86.1\%$ isotropic --- a scalar times
the identity --- where a Newtonian tidal tensor is traceless. Separately, no quadrupole contribution to
the tidal tensor was implemented anywhere, so the ``tidal'' far field was monopole-only even where the
force used the quadrupole. Consequently the earlier statement that the tidal error is $\approx2\times$
the force error was an artefact: the corrected factor is $5$--$6\times$, and the leading error's
$\theta$-power is $3.56$ (force: $3.23$), so the ``anomalous'' second-derivative $\theta$-scaling
claimed in an earlier draft does not exist. One canonical reference now exists
(\texttt{src/\allowbreak laniakea\_\allowbreak hmmp/\allowbreak models/\allowbreak exact\_\allowbreak field.\allowbreak py}), pinned to the autograd Jacobian by eleven tests,
and the operator carries both derivatives including the quadrupole term. Force numbers are unaffected:
the re-run reproduced every published force error bit-for-bit.

\noindent\textbf{(D7) The novelty claim is scoped to the constants and the geometry, not to the
existence of a remainder.} An earlier draft of this work claimed that no second-derivative remainder
exists, on the grounds that we could not find one. \textbf{That claim is withdrawn, because
Appendix~\ref{app:secondorder} derives one.} What survives is narrower and defensible. The classical
FMM error-analysis line --- \citet{cheng1999fast}, and the 3-D potential
convergence theorems that follow it --- is the most plausible place such a derivative bound already
exists. \textbf{We have now read the 1999 paper, and it does not contain one: it states the classical
\emph{potential} remainder only.} The five remainder statements, their equation numbers, the basis
difference and the three reasons they do not differentiate into ours are set out at the head of
Appendix~\ref{app:secondorder}, which is where the comparison belongs; the exhaustive probe counts are
in \texttt{docs/\allowbreak PROVENANCE.\allowbreak md} \S4d. What remains unread is the literature that \emph{follows} that
line rather than the line's origin.

The five further candidates that an earlier access failure had left unread have also since been
obtained and read, including the one that mattered most: a \emph{learned} multipole
translation operator, which is where a truncation bound on learned hierarchical machinery would have
lived. \textbf{All five report measured accuracy and state no bound}; none bounds a derivative, and
none carries an order-dependent constant. We therefore claim novelty of the
\emph{explicit per-order
constants}, of the centre-of-mass-versus-geometric-size distinction, and of the per-interaction
verification against measured remainders --- \textbf{not} priority for the existence of a derivative
remainder. Should such a bound exist in work we have not read, the contributions that stand are still
the law's explicit constants and geometry content, the measurement at $2.3$M, and the demonstration
that the measured order-shift product is data-dependent (Section~\ref{sec:order}).

\noindent\textbf{(D8) Our octree's recorded cell centres were offset by a full cell side.} The
recorded centre of every non-root cell --- the point both geometric acceptance tests measure their
acceptance distance to --- was wrong. The subdivision stepped from a parent centre by $s/2$ where an
octant centre sits at $s/4$, so the recorded centre landed a full cell side from the box whose
particles it describes, and the error compounds down the tree. Measured on three point sets
(uniform random, Gaussian, and a two-body-dominated scene), \textbf{$97$--$99\%$ of cells had member
particles outside $[\text{centre}\pm\text{size}/2]$}, with a worst-case offset of exactly $2.0$
half-sides. Two further consequences followed: the octant \emph{partition} below depth one also used
the wrong centre, so the tree was not a spatial subdivision; and callers centred the root on the
positional \emph{mean} while sizing it by the maximum \emph{span}, so on a non-uniform cloud
particles escaped the root box itself. The multipole \emph{moments} were computed from particle
membership and were never affected --- which is why the published centre-of-mass test, added late,
looked healthy while the geometric tests did not.

\emph{How it surfaced, and what it cost us.} The symptom was a catastrophic, intermittent blow-up at
$\theta=0.70$ in the geometric-centre test, which an earlier draft of this paper reported as a
limitation of the Barnes--Hut acceptance criterion. That reading was wrong: the published criterion
measures the distance to the cell's centre of mass, and none of the production tree codes we checked
(GADGET-2/3/4, SWIFT, AREPO, GIZMO, PKDGRAV3, ChaNGa) uses a geometric-centre acceptance test, so
none of them is exposed. We withdrew the failure-mode contribution in full. \textbf{The correction is
worth $11$--$33\times$ in accuracy at fixed $\theta$ and costs nothing}: at $N=2{,}000$, $\theta=0.5$
the global relative force error falls from $0.021419$ to $0.001996$; at $N=10^{5}$ the per-node force
and tidal errors fall by $11\times$ and $15\times$ ($p{=}0$) and $22\times$ and $33\times$ ($p{=}1$).

\emph{Why it survived every existing test.} A two-body orbit keeps both particles in one leaf; a
shallow tree admits few far cells; and the moment aggregation reads particle membership rather than
cell geometry, so nothing in the suite ever compared a cell centre against its contents. The
invariant that catches it is independent of any multipole approximation --- every cell's particles
must lie inside $[\text{centre}\pm\text{size}/2]$ --- and is now asserted by
\texttt{tests/\allowbreak test\_\allowbreak cell\_\allowbreak geometry.\allowbreak py}, which also asserts that the pre-fix construction still
\emph{reproduces} the defect so that artifacts predating the fix stay interpretable. Every
\emph{accuracy} number in this paper is post-D8; \texttt{legacy\_geometry=True} restores the old
behaviour. \textbf{The two cost tables are the exception}: the peak-memory and forward-pass-time
series resolve to artifacts that carry no post-D8 marker and are not part of the released set
(\texttt{outputs/\allowbreak memory\_\allowbreak hmmp\_\allowbreak 2M3\_\allowbreak *.\allowbreak json}, \texttt{outputs/\allowbreak scaling.\allowbreak json}), so their post-fix
status is not established here. The geometry correction changes the tree, so both will move.
\texttt{docs/\allowbreak PROVENANCE.\allowbreak md} \S3.5 tracks them.

\subsection{Scope of the D8 correction: what is post-fix, and what is not}
\label{app:d8scope}

The geometry fix changes the tree, and therefore every measured error and every convergence curve.
It does \emph{not} change the algebra of the second-derivative bound or its exact per-order
constants, which are properties of the multipole expansion. Table~\ref{tab:d8scope} separates the two.

\begin{table}[h]
\centering
\small
\begin{tabular}{p{0.50\columnwidth}p{0.40\columnwidth}}
\toprule
claim / result & status with respect to D8 \\
\midrule
bound, exact constants, per-interaction verification & geometry-independent; verified without the tree \\
all accuracy/error numbers cited here (order shift, quadrupole, eigen/flip, gradient, 16.7M) & post-D8; re-run or recomputed with the corrected tree \\
the D8 correction itself ($11$--$33\times$) & post-fix minus pre-fix; the only pre-fix numbers are quoted as that comparison \\
two cost series (forward time, peak memory) & \textbf{not established post-D8}; excluded from the claims and flagged in the text \\
periodic-boundary statement & convention, independent of D8 \\
\bottomrule
\end{tabular}
\caption{Which claims depend on the corrected octree geometry. The bound's validity is independent
of D8; the numerical error levels and convergence curves are not.}
\label{tab:d8scope}
\end{table}

\noindent\textbf{Summary.} Every \emph{accuracy} number cited in this paper is
post-D8. The fix moves the convergence curves and the absolute errors; it does not invalidate the
second-derivative bound, whose constants are exact and whose verification is per interaction. The
only series not established post-D8 are the two cost tables, which the text already marks as such.

\section{Comparing acceptance criteria after the geometry fix}\label{app:surface}
Three acceptance tests are implemented: distance to the cell's centre of mass (\texttt{com}, the
published criterion), distance to the geometric cell centre (\texttt{geo}), and distance to the
cell's circumsphere boundary (\texttt{surface}). Section~\ref{sec:limits} explains why an earlier
draft measured the second of these, mistook defect D8 for a property of Barnes--Hut, and reported
the resulting blow-up as a limitation. With the geometry corrected the three can be compared as an
accuracy-versus-cost question rather than as a failure and a fix. Two-body-dominated scene,
$N=600$, $Q=256$, $p{=}0$, same particles, same tree and same operator, only the acceptance test
changed:

\begin{center}
\resizebox{\ifdim\width>\linewidth\linewidth\else\width\fi}{!}{%
\begin{tabular}{lrrrr}
\toprule
& \multicolumn{2}{c}{force rel.\ err} & \multicolumn{2}{c}{tidal rel.\ err} \\
$\theta$ & 0.50 & 0.70 & 0.50 & 0.70 \\
\midrule
com     & $1.08\times10^{-4}$ & $2.96\times10^{-4}$ & $2.92\times10^{-3}$ & $8.01\times10^{-3}$ \\
geo     & $1.32\times10^{-4}$ & $4.70\times10^{-4}$ & $3.26\times10^{-3}$ & $9.28\times10^{-3}$ \\
surface & $\mathbf{3.73\times10^{-5}}$ & $\mathbf{8.38\times10^{-5}}$ & $\mathbf{1.24\times10^{-3}}$ & $\mathbf{2.26\times10^{-3}}$ \\
\bottomrule
\end{tabular}
}
\end{center}

\noindent Two things follow. \textbf{(i) The published criterion is the better of the two single-point
tests, and the circumsphere test is better still} --- by $2.9\times$ on the force and $2.4\times$ on
the tidal tensor at $\theta=0.50$ relative to \texttt{com}, at the cost of more far-field entries per
query node. We report that as a trade, not as a correctness fix. \textbf{(ii) None of the three fails
at any $\theta$ we tested}, $\theta=0.70$ included; the catastrophe earlier drafts attributed to the
acceptance test was defect D8 and is withdrawn.

The \texttt{surface} test is \emph{not} novel: an analytically bounded multipole acceptance criterion
is due to Warren \& Salmon (1993), Dehn\-en (2002) advertises a simplified such criterion, and ChaNGa
ships a sphere--box surface test. Our contribution here is the measurement under matched conditions,
not the criterion.

These numbers regenerate with \texttt{scripts/\allowbreak \_\allowbreak ab\_\allowbreak cell\_\allowbreak centres.\allowbreak py}, which reports all three
criteria on both geometries and prints the pre-fix comparison alongside.

\noindent\textbf{Table~\ref{tab:eigen}: the physical metrics behind Section~\ref{sec:physics}.}
Principal-axis misalignment is quoted only over nodes whose top two $|\lambda|$ differ by at least
$1.5\times$ ($29$--$40\%$ of nodes here), because an eigenvector is meaningless where the eigenvalues
are degenerate; the unrestricted p95 is larger by about half again (e.g.\ $0.68^{\circ}$ against
$0.41^{\circ}$ for the $N=2.3$M, $\theta=0.5$, $p{=}0$ com row), which is itself a reminder that a
bare angle is a trap. The
exact field is $71\%$ sheet and $29\%$ filament, so the flip column is a classification change, not a
boundary effect.

\begin{table}[h]
\centering
\caption{Eigenvalue error and principal-axis misalignment (p95, well-separated nodes only) of the
approximate tidal tensor against the exact one, on the real CAMELS snapshot, $Q=4096$, seed $0$,
recomputed with the post-defect-D8 tree: \texttt{scripts/\allowbreak eval\_\allowbreak tidal\_\allowbreak eigenmetrics.\allowbreak py}. ``com'' is the
centre-of-mass acceptance test, ``cell'' the distance-to-cell test. \textbf{The two criteria now agree
to a few per cent, and neither fails}: at $\theta=0.70$ --- well past the classical safety line
$\theta<1/\sqrt{3}$ --- the worst misalignment anywhere in the table is $1.3^{\circ}$. An earlier
revision carried pre-D8 numbers under a \texttt{com} label that the generating script could not then
produce, because its \texttt{--opening-mode} choices omitted \texttt{com} and defaulted to
\texttt{center}; those rows are replaced.}\label{tab:eigen}
\resizebox{\ifdim\width>\linewidth\linewidth\else\width\fi}{!}{%
\begin{tabular}{lrrrrr}
\toprule
$N$ & $\theta$ & opening & $p$ & eig.\ rel.\ err & p95 axis (deg) \\
\midrule
$2.3$M & 0.5 & com & 0 & $3.79\times10^{-3}$ & 0.41 \\
$2.3$M & 0.5 & com & 1 & $1.19\times10^{-3}$ & 0.13 \\
$100$k & 0.3 & com & 0 & $5.01\times10^{-4}$ & 0.05 \\
$100$k & 0.5 & com & 0 & $3.13\times10^{-3}$ & 0.39 \\
$100$k & 0.5 & com & 1 & $1.02\times10^{-3}$ & 0.12 \\
$100$k & 0.5 & cell & 0 & $3.23\times10^{-3}$ & 0.38 \\
$100$k & 0.7 & com & 0 & $1.08\times10^{-2}$ & 1.22 \\
$100$k & 0.7 & cell & 0 & $1.17\times10^{-2}$ & 1.31 \\
$100$k & 0.7 & cell & 1 & $5.66\times10^{-3}$ & 0.67 \\
\bottomrule
\end{tabular}
}
\end{table}

\begin{table}[h]
\centering
\caption{Classification fidelity, the column that matters to a user of the tidal tensor: the fraction
of query nodes whose eigenvalue-sign class (the number of positive eigenvalues, which for a traceless
tensor is $1$ or $2$ --- see Section~\ref{sec:physics}) differs from the exact one, for the same rows
as Table~\ref{tab:eigen}. Post-D8 the worst rate in the table is $1.4\%$, and at fixed
$\theta=0.70$ the acceptance test no longer separates the two criteria --- what used to look like a
criterion effect was defect~D8.}\label{tab:flip}
\resizebox{\ifdim\width>\linewidth\linewidth\else\width\fi}{!}{%
\begin{tabular}{lrrrr}
\toprule
$N$ & $\theta$ & opening & $p$ & class flip \\
\midrule
$2.3$M & 0.5 & com & 0 & $0.51\%$ \\
$2.3$M & 0.5 & com & 1 & $0.27\%$ \\
$100$k & 0.3 & com & 0 & $0.10\%$ \\
$100$k & 0.5 & com & 0 & $0.17\%$ \\
$100$k & 0.5 & com & 1 & $0.15\%$ \\
$100$k & 0.5 & cell & 0 & $0.34\%$ \\
$100$k & 0.7 & com & 0 & $1.39\%$ \\
$100$k & 0.7 & cell & 0 & $1.39\%$ \\
$100$k & 0.7 & cell & 1 & $0.71\%$ \\
\bottomrule
\end{tabular}
}
\end{table}

\section{The acceptance criterion on dynamically hard scenes}\label{app:scenes}
Appendix~\ref{app:surface} compares the three acceptance tests on one scene at $\theta\le0.70$.
A centre-of-mass test fails when a cell's mass is poorly separated from the query particle, and that
is a property of the \emph{scene}, not of $\theta$ alone --- so the comparison is repeated on five,
sweeping $\theta$ well past the classical $\theta<1/\sqrt{3}\approx0.577$ safety line. Scenes: a
lognormal cloud, plus four of the sibling project's physical presets, which have deep unequal wells
and close encounters. $N=600$ ($20{,}000$ for the lognormal), $Q=256$, three query seeds, $p{=}0$,
against the brute-force reference.

\begin{center}
\resizebox{\ifdim\width>\linewidth\linewidth\else\width\fi}{!}{%
\begin{tabular}{lrrrr}
\toprule
scene & heaviest-particle mass share & \multicolumn{3}{c}{first $\theta$ at which the test fails} \\
 & & com & geometric centre & circumsphere \\
\midrule
synthetic lognormal   & $\sim0$ & none $\le1.20$ & none $\le1.20$ & none $\le1.20$ \\
galaxy\_collision     & $7.1\%$  & none $\le1.20$ & none $\le1.20$ & none $\le1.20$ \\
spiral\_galaxy        & $13.6\%$ & none $\le1.20$ & $\mathbf{1.20}$ & none $\le1.20$ \\
binary\_star          & $30.1\%$ & none $\le1.20$ & $\mathbf{1.20}$ & none $\le1.20$ \\
globular\_cluster     & $45.5\%$ & $0.90$ & none $\le1.20$ & none $\le1.20$ \\
\bottomrule
\end{tabular}
}
\end{center}

\noindent Three readings.

(i) \textbf{At the settings users run, there is no failure to report.} The published \texttt{com}
criterion is clean on every scene at $\theta\le0.80$, and production tree codes default to
$\theta\approx0.7$ (all of them \texttt{com}-anchored, several using a relative force-accuracy
criterion instead of a fixed $\theta$ at all). This is the honest scope of the failure-mode question:
the $\theta=0.70$ collapse reported in earlier drafts was defect D8
(Appendix~\ref{app:provenance}), and with the geometry corrected nothing takes its place in the
production range.

(ii) \textbf{The geometric-centre test does have a real boundary, and it is at the analytically
predicted place.} For that test the minimum separation between the query particle and an admitted
cell's \emph{expansion} centre is $\text{size}\cdot(1/\theta-\sqrt{3}/2)$, which vanishes at
$\theta=2/\sqrt{3}\approx1.155$. Above it, admission no longer implies the particle is outside the
source's bounding sphere. The measured onset is $\theta=1.20$ on the two scenes with the sharpest
central concentration --- consistent with the prediction, and it appears on exactly the scenes the
mechanism requires. \textbf{The \texttt{com} test has a tighter line rather than no line:} admission
gives $r>\text{size}/\theta$ measured to the \emph{centre of mass}, and because the centre of mass may
sit anywhere inside the box the cell's contents can reach $\sqrt3\,\text{size}$ from it, so strict
outside-ness is guaranteed only for $\theta\le1/\sqrt3\approx0.577$ --- the classical line, half the
geometric-centre bound (Section~\ref{sec:criteria} derives both). \textbf{So the two single-point tests
are not equivalent, and the published one carries the tighter guarantee}, which is the sense in which
it is the safe one; but since no production code uses a
geometric-centre acceptance test, this bounds a hypothetical rather than a practice, and it is why we
report it here rather than in the main text.

(iii) \textbf{The circumsphere test is safe everywhere and the most accurate.} Across every scene and
every $\theta$ up to $1.20$ its direction cosine stays $\ge0.99998$, and its error is $3$--$10\times$
below \texttt{com}'s (Appendix~\ref{app:surface}).

\noindent\textbf{One anomaly, traced and resolved.} On \texttt{globular\_cluster} the published
\texttt{com} criterion shows a sharp jump between $\theta=0.80$ and $0.90$, and it is
\emph{seed-dependent}: at $\theta=0.90$ the force error is $1.2\times10^{-2}$ with cosine $0.99992$
for one query draw and $4.7\times10^{-1}$ with cosine $0.88411$ for another, on the same particles and
the same tree. It is not a criterion failure, and two measurements settle it. First, the error is not
concentrated: the worst single node carries $5.3\%$ of the total per-node error, and the per-node
\emph{median} moves smoothly across the onset ($2.74\times10^{-2}\to3.36\times10^{-2}$) --- a genuine
blow-up is one node dominating, as the D8 artifact was. Second, the admitted cells are legitimate: the
closest has $|x_{\mathrm{com}}-x|=0.454$ against $\text{size}=0.395$, i.e. $\text{size}/d=0.871<0.90$,
admitted by design, with the query particle outside the cell. Its circumsphere radius is $0.342$, so
the particle sits $0.11$ radii outside the source, where the monopole error is
$O(\text{size}/d)\approx0.87$. That is what $\theta=0.90$ \emph{means} for a monopole expansion.

The seed dependence is then a measurement artifact, and a methodological one: the \emph{global
direction cosine} over $Q=256$ sampled nodes is fragile --- whether the sample happens to include
nodes in the near/far transition region swings it from $0.99992$ to $0.884$, while the per-node median
barely moves. The query-sampling study of Appendix~\ref{app:repro} quotes $1.2$--$1.7\%$ on the
per-node mean; the cosine is far more sensitive, and should not be quoted from a small sample. This is
the same lesson as the norm-weighted-statistics warning of Section~\ref{sec:limits}, in a second
place.

A controlled sweep over heaviest-particle mass fraction ($f=0.005$ to $0.60$, heavy particle at the
centre) reproduced \emph{no} onset at any $f$, rejecting the obvious hypothesis that concentration
alone governs it: with the heavy mass at the centre of a symmetric cloud, the expansion centre of
every enclosing cell sits near the centre and ${\rm size}/d$ is large, so those cells are refused.
What matters is a cell whose expansion centre lands close to a query particle while the cell stays
large, which needs an off-centre concentration rather than merely a concentrated one.

\noindent Artifact: \texttt{outputs/\allowbreak d2\_\allowbreak scene\_\allowbreak threshold\_\allowbreak sweep.\allowbreak json};
scripts \texttt{scripts/\allowbreak d2\_\allowbreak scene\_\allowbreak threshold\_\allowbreak sweep.\allowbreak py},
\texttt{scripts/\allowbreak d2\_\allowbreak mass\_\allowbreak concentration\_\allowbreak threshold.\allowbreak py}, \texttt{scripts/\allowbreak \_\allowbreak d2\_\allowbreak trace\_\allowbreak globular.\allowbreak py}.

\section{Does the second-derivative error land where it matters?}\label{app:environment}
The class-flip rate of Section~\ref{sec:physics} is a bare number until it is weighted by environment.
A misclassified node in an empty region costs nothing; one on a filament spine is a statement about
where galaxies form. We therefore bin the same query nodes by a $k$-NN density proxy ($k=32$,
$\rho\propto r_k^{-3}$ with the self excluded) and report the flip rate, the share of all flips, the
eigenvalue-sign mixture and the eigenvalue scale per octile. Real CAMELS snapshot, $N=100{,}000$,
$Q=4{,}096$, $\theta=0.5$, centre-of-mass criterion, $p{=}0$; artifact
\texttt{outputs/\allowbreak misclassification\_\allowbreak vs\_\allowbreak environment.\allowbreak json}, script
\texttt{scripts/\allowbreak eval\_\allowbreak misclassification\_\allowbreak vs\_\allowbreak environment.\allowbreak py}.

\begin{center}
\resizebox{\ifdim\width>\linewidth\linewidth\else\width\fi}{!}{%
\begin{tabular}{rrrrrr}
\toprule
density octile & median $\rho$ & exact mix (1-pos/2-pos) & flip rate $p{=}0$ & flip rate $p{=}1$ & share of flips \\
\midrule
0 (lowest) & 14.0 & $0.77/0.23$ & $0.0645$ & $0.0488$ & $0.121$ \\
1 & 18.5 & $0.77/0.23$ & $0.0879$ & $0.0625$ & $0.165$ \\
2 & 22.4 & $0.76/0.24$ & $0.0723$ & $0.0469$ & $0.136$ \\
3 & 26.5 & $0.80/0.20$ & $0.0742$ & $0.0508$ & $0.139$ \\
4 & 31.6 & $0.77/0.23$ & $0.0508$ & $0.0371$ & $0.095$ \\
5 & 38.7 & $0.76/0.24$ & $0.0684$ & $0.0352$ & $0.128$ \\
6 & 48.7 & $0.82/0.18$ & $0.0645$ & $0.0449$ & $0.121$ \\
7 (highest) & 77.3 & $0.79/0.21$ & $0.0508$ & $0.0352$ & $0.095$ \\
\bottomrule
\end{tabular}
}
\end{center}

\noindent The flip rate spans only $5.1$--$8.8\%$ across a factor $5.5$ in density, with no monotone
trend, and the class mixture is flat at $76$--$82\%$ sheet-like in every octile. The low-density half
carries $0.140$ of all flips against the high-density half's $0.110$ (a uniform distribution would
give $0.125$ each), so there is at most a weak tilt towards low density and it is not the dominant
effect. \textbf{We therefore report no evidence that the second-derivative error concentrates in the
environments that matter}, and equally no evidence that it avoids them. Two caveats keep this from
being a stronger statement: the snapshot is at $z\approx6$, where the density contrast is mild, and
the $k$-NN density on a $100$k subsample is a noisy environment proxy. A cluster-scale or low-redshift
snapshot would be the right place to repeat it.

\section{The neighbour-cutoff baseline}\label{app:baseline}
Real CAMELS snapshot, $N=100{,}000$, $Q=2{,}048$ query nodes for the accuracy columns and
$N=20{,}000$ over \emph{all} nodes for the momentum column, $\theta=0.5$, $\texttt{max\_leaf}=8$,
post-D8 geometry. The $k$-nn baseline sums the $k$ nearest neighbours of each node directly; $k$ is
set to the operator's measured terms per query node ($664.6$, i.e.\ $1{,}145{,}335$ far entries plus
$215{,}667$ near-leaf member particles over $2{,}048$ nodes) so the two cost the same number of
evaluations. Script \texttt{scripts/\allowbreak eval\_\allowbreak knn\_\allowbreak baseline.\allowbreak py}, artifact
\texttt{outputs/\allowbreak knn\_\allowbreak baseline\_\allowbreak postd8.\allowbreak json}.

\begin{center}
\resizebox{\ifdim\width>\linewidth\linewidth\else\width\fi}{!}{%
\begin{tabular}{lrrrrr}
\toprule
engine & terms/node & force rel.\ err & tidal rel.\ err & tidal/force & momentum violation \\
\midrule
tree, $p{=}0$ & $664.6$ & $3.02\times10^{-3}$ & $9.97\times10^{-3}$ & $3.31$ & $9.2\times10^{-5}$ \\
tree, $p{=}1$ & $664.6$ & $5.34\times10^{-4}$ & $2.94\times10^{-3}$ & $5.51$ & $9.6\times10^{-6}$ \\
$k$-nn cutoff, $k{=}665$ & $665$ & $9.32\times10^{-1}$ & $1.79\times10^{-1}$ & $0.19$ & $4.5\times10^{-3}$ \\
\midrule
exact sum, for reference & --- & $0$ & $0$ & --- & $3.6\times10^{-8}$ \\
\bottomrule
\end{tabular}
}
\end{center}

\noindent The momentum column is $\|\sum_i m_i a_i\| / \sum_i m_i\|a_i\|$ from a single force
evaluation: the exact sum conserves it to round-off, the tree violates it structurally (its far field
omits the reaction), and the cutoff violates it \emph{more} --- $49\times$ more at $p{=}0$ --- because
truncation is not antisymmetric. See Section~\ref{sec:conservation} for what that violation does over
an integration. \emph{These numbers, and those of Appendix~\ref{app:critic}, are post-D8; the
corresponding pre-D8 artifacts are retained as \texttt{knn\_\allowbreak baseline.\allowbreak json} and
\texttt{lumyn\_\allowbreak critic\_\allowbreak bridge.\allowbreak json} for comparison, and differ by about an order of magnitude ---
always in the direction of the corrected geometry being better.}


\noindent\textbf{The mesh route, which is the one cosmology actually uses.} A $k$-NN cutoff is not how
anyone computes gravity. The standard cheap solver is \emph{particle--mesh} (PM): deposit the masses on
a grid, solve the Poisson equation with FFTs, difference the potential for the force, and interpolate
back. It is $\mathcal{O}(N + M^3\log M)$ and it is what every production cosmology code uses for its
long-range force, so the fair question is not whether a cutoff competes but whether a mesh does.
We implemented one (\texttt{scripts/\allowbreak eval\_\allowbreak pm\_\allowbreak baseline.\allowbreak py}) and validated the solver independently
before trusting any comparison: on a smooth analytic density $\rho\propto\cos(kx)$ it reproduces the
exact force to $2\times10^{-4}$ at $k\,dx=0.10$ and $1.3\times10^{-2}$ at $k\,dx=0.79$, i.e.\ with the
$(k\,dx)^2$ signature a second-order finite-difference Poisson solver must have
(\texttt{scripts/\allowbreak \_\allowbreak diag\_\allowbreak pm\_\allowbreak smooth.\allowbreak py}).

\textbf{Two conventions had to be reconciled first, and getting this wrong inflates PM's error by
$80$ percentage points.} A periodic FFT-Poisson solve drops the $k{=}0$ mode, which is the field of
the box's own mean density; an isolated tree and a brute-force pair sum both carry it. In this box
($L=25$ Mpc/$h$, $\bar\rho=0.955$) that term has mean magnitude $44.6$, against an isolated force of
$54.6$ --- so it is $2.3\times$ the \emph{peculiar} part and dominates the total. Measuring a periodic
PM against an isolated reference therefore reports $98\%$ where the mesh is only responsible for
$18\%$. We compute the uniform-box field by summation rather than from Gauss's law, because
$(4\pi G\bar\rho/3)\,r$ is the field inside a uniform \emph{sphere} and this box is a cube, where that
formula is $8\%$ wrong (\texttt{scripts/\allowbreak \_\allowbreak diag\_\allowbreak pm\_\allowbreak peculiar.\allowbreak py}). The table below puts PM on the
isolated convention, the same one the operator is measured in.

\begin{center}
\resizebox{\ifdim\width>\linewidth\linewidth\else\width\fi}{!}{%
\begin{tabular}{lrrrrr}
\toprule
engine & part./cell & force rel.\ err & tidal rel.\ err & cost (s) & what the cost buys \\
\midrule
tree, $p{=}0$, $\theta{=}0.5$ & --- & $\mathbf{2.76\times10^{-3}}$ & $\mathbf{3.19\times10^{-2}}$ & $31.4$ & tree + $1{,}024$-node field \\
PM $64^{3}$ & $8.77$ & $1.98\times10^{-1}$ & $9.60\times10^{-1}$ & $6.4$ & the whole $2.3$M-node field \\
PM $128^{3}$ & $1.10$ & $1.90\times10^{-1}$ & $9.12\times10^{-1}$ & $8.8$ & the whole $2.3$M-node field \\
PM $256^{3}$ & $0.14$ & $1.80\times10^{-1}$ & $7.84\times10^{-1}$ & $21.5$ & the whole $2.3$M-node field \\
\bottomrule
\end{tabular}
}
\end{center}

\noindent Real CAMELS snapshot, $N=2.3$M, $Q=1{,}024$, $\theta=0.5$, $\texttt{max\_leaf}=8$, per-node
mean relative error, $\epsilon=10^{-4}$; artifacts \texttt{outputs/\allowbreak pm\_\allowbreak baseline\_\allowbreak n2M3.\allowbreak json} and, at
$N=10^{5}$, \texttt{outputs/\allowbreak pm\_\allowbreak baseline\_\allowbreak n100k.\allowbreak json}. The tree row costs more than any PM row while
evaluating $1{,}024$ nodes rather than $2.3$M: \textbf{the mesh is roughly two orders of magnitude
cheaper, and that is the honest half of this appendix.}

\noindent\textbf{What PM's error is, and is not.} Refining the mesh by $16\times$ moves the force error
from $1.98$ to $1.80\times10^{-1}$: it does not converge. That is not a defect of our implementation
--- the analytic control above rules that out --- it is what a mesh method does to this quantity. A
mesh represents the \emph{field}; the per-node acceleration in a clustered particle distribution is
dominated by structure \emph{below the cell}, which a deposited density cannot carry and a `cloud'
assignment cannot read back. Against a reference smoothed at the cell scale PM's error improves but
never collapses (best $53\%$ over $\epsilon=10^{-4}$ to $1.6$;
\texttt{scripts/\allowbreak \_\allowbreak diag\_\allowbreak pm\_\allowbreak smoothing.\allowbreak py}), and the residual against an unsmoothed reference is
exactly the sub-cell part that production codes supply with a short-range direct sum. That pairing is
tree--PM, and its accuracy at small scales is set by the short-range cutoff --- which is the regime the
$k$-NN baseline above already characterises.

\noindent\textbf{The tidal tensor is where a mesh is worst.} PM's tidal error is $4$--$5\times$ its
force error at every mesh ($0.78$--$0.96$ against $0.18$--$0.20$), and again barely responds to
resolution. The tidal tensor is a \emph{difference} of forces, so it weights the small-scale structure
that the deposition removes more heavily than the force does. This is the same ordering the operator
itself shows at scale (its tidal error is $2.6$--$10.8\times$ its force error), and it is the reason a
mesh is a poor substitute for this operator's contract even where it is affordable.

\noindent\textbf{One structural difference worth recording.} PM's tidal tensor is not traceless: its
trace is $4\pi G(\rho_{\rm dep}-\bar\rho)$, which the finite-difference construction returns exactly
and which at coarse particle loading is dominated by deposition shot noise (mean $|\mathrm{tr}|$ rises
from $8.3$ at $M=64$ to $34.9$ at $M=256$ as particles per cell fall from $8.8$ to $0.14$, tracking
$4\pi G\bar\rho/\sqrt{N_{\rm pc}}$). The brute-force pair sum, and therefore this operator, is traceless
to machine precision because it is the Jacobian of a softened pairwise kernel. Neither is wrong; they
answer different questions, and the trace-removed PM column above is the comparison that does not
depend on which convention is chosen.

\section{Cross-check with an independent conservation critic}\label{app:critic}
Section~\ref{sec:conservation} concludes that the operator's non-conservativity is a momentum effect
and not an energy effect. That conclusion was reached from the solver side; here it is checked from
the critic side, using an independent implementation: a conservation-law critic built for a
physically-grounded scene-generation engine, whose four normalised drift channels are energy,
angular momentum, momentum and centre of mass. The critic's default weights are angular $+$ momentum,
chosen because its authors measured angular and momentum to be $88$--$100\%$ and $85$--$100\%$
discriminative while \emph{energy was only $56$--$60\%$}. Feeding it the trajectories of
Section~\ref{sec:conservation} (same initial conditions, $150$ steps, $\theta=0.5$, softening
$0.05$, $N=2{,}000$) reproduces their finding from the other end:

\begin{center}
\resizebox{\ifdim\width>\linewidth\linewidth\else\width\fi}{!}{%
\begin{tabular}{lrrrr}
\toprule
force used & energy & angular & momentum & centre of mass \\
\midrule
exact & $1.69\times10^{-4}$ & $1.20\times10^{-6}$ & $1.94\times10^{-8}$ & $1.30\times10^{-8}$ \\
operator, $p{=}0$ & $6.04\times10^{-4}$ & $2.33\times10^{-3}$ & $1.90\times10^{-4}$ & $9.81\times10^{-6}$ \\
operator, $p{=}1$ & $5.12\times10^{-4}$ & $3.83\times10^{-4}$ & $1.76\times10^{-5}$ & $1.31\times10^{-6}$ \\
\midrule
ratio to exact, $p{=}0$ & $3.6$ & $1.9\times10^{3}$ & $9.8\times10^{3}$ & $7.6\times10^{2}$ \\
\bottomrule
\end{tabular}
}
\end{center}

\noindent Swapping the exact force for this operator moves the energy channel by a factor of $3.6$
and the two channels that critic weights by three to four orders of magnitude. So the channel weighting
is not a heuristic: \emph{it follows from the structure of the force approximation}. One detail is
worth recording because it was not obvious to us: \textbf{the angular-momentum channel is the more
sensitive of the two} ($12\times$ the momentum drift here, $20\times$ under the pre-D8 geometry), since
the reaction force this operator omits would act on a cell's members at large distance, so the torque
error carries an extra factor of $r$. Both channels fall by $6$--$11\times$ when the quadrupole is
retained, consistent with $p{=}1$ reducing the far-field error throughout. Script
\texttt{scripts/\allowbreak bridge\_\allowbreak lumyn\_\allowbreak critic.\allowbreak py}, artifact
\texttt{outputs/\allowbreak lumyn\_\allowbreak critic\_\allowbreak bridge\_\allowbreak postd8.\allowbreak json}.

\noindent\textbf{The solver-side measurement.} The table summarised in
Section~\ref{sec:conservation} is the over-time counterpart of the critic channels above: $150$
leapfrog steps at $N=2{,}000$, $\theta=0.5$, softening $0.05$, initial conditions with zero net
momentum (\texttt{scripts/\allowbreak eval\_\allowbreak conservation\_\allowbreak drift.\allowbreak py}, \texttt{outputs/\allowbreak conservation\_\allowbreak drift.\allowbreak json}).

\begin{center}
\resizebox{\ifdim\width>\linewidth\linewidth\else\width\fi}{!}{%
\begin{tabular}{lrrr}
\toprule
force used & $\|\sum_i m_i v_i\|$ & energy drift $|\Delta E/E_0|$ & barycentre $\|\bar x\|$ \\
\midrule
exact & $1.2\times10^{-4}$ & $5.0\times10^{-5}$ & $1.1\times10^{-6}$ \\
operator, $p{=}0$ & $1.4\times10^{0}$ & $2.3\times10^{-4}$ & $4.6\times10^{-4}$ \\
operator, $p{=}1$ & $1.1\times10^{-1}$ & $1.6\times10^{-4}$ & $6.2\times10^{-5}$ \\
\bottomrule
\end{tabular}
}
\end{center}

\noindent Momentum is violated by $1.1\times10^{4}\times$ the exact run's, and the two series are
qualitatively different rather than merely different in size: the operator's violation is monotone and
linear in time, whereas the exact run's is bounded and does not grow --- it is the cancellation floor of
a near-zero vector sum in \texttt{float32}, the integrator being momentum-symmetric. It is a
structural defect that more multipole order reduces ($12\times$ at $p{=}1$) but cannot remove, while
energy drift stays the \emph{same order} as the exact run's ($4.7\times$ at $p{=}0$, $3.3\times$ at
$p{=}1$), because it is set by the integrator and the approximate force remains close enough to the
gradient of some potential that no secular energy error appears. Script
\texttt{scripts/\allowbreak eval\_\allowbreak conservation\_\allowbreak drift.\allowbreak py}, artifact
\texttt{outputs/\allowbreak conservation\_\allowbreak drift\_\allowbreak postd8.\allowbreak json}; the pre-D8 counterpart
\texttt{conservation\_\allowbreak drift.\allowbreak json} gave a $p{=}0$ momentum violation of $10.4$, and is retained for
comparison.

\section{Reproducibility details}\label{app:repro}
\noindent\textbf{Data.} CAMELS IllustrisTNG $L25n256$, realisation \texttt{CV\_14},
\texttt{snapshot\_014} ($z\approx6$, box $25$\,Mpc/$h$, $\Omega_m=0.3$, $h=0.6711$). Download:
\texttt{scripts/\allowbreak download\_\allowbreak camels.\allowbreak py} (public Flatiron mirror; the CAMELS public data release
\citep{camels2023} is the durable route if the mirror path changes). Reader/adapter:
\texttt{src/\allowbreak laniakea\_\allowbreak hmmp/\allowbreak data/\allowbreak camels.\allowbreak py}. $\mathrm{PartType1}$ dark matter, $16.7$M particles
subsampled to $2.3$M.

\noindent\textbf{Operator configuration.} Opening angle $\theta=0.5$ and monopole only ($p=0$) unless
a setting is explicitly varied; softening $\epsilon=10^{-4}$; octree \texttt{max\_leaf}$=16$,
\texttt{max\_depth}$=20$ (sweeps) or $24$ (exact-N-body check). Positions rescaled to
$[0,25)$\,Mpc/$h$. The octree and interaction list are built once per particle set and cached across
repeated evaluations, which is the configuration the memory numbers assume.

\begin{sloppypar}
\noindent\textbf{Script $\rightarrow$ result.} \texttt{scripts/\allowbreak exact\_\allowbreak nbody.\allowbreak py}
(Table~\ref{tab:exact}); \texttt{scripts/\allowbreak sweep\_\allowbreak multipole\_\allowbreak order.\allowbreak py}
(Table~\ref{tab:ordersweep} and the $\theta=0.70$ breakdown);
\texttt{scripts/\allowbreak exact\_\allowbreak nbody\_\allowbreak sampled.\allowbreak py} via \texttt{scripts/\allowbreak rerun\_\allowbreak sampled\_\allowbreak validation.\allowbreak sh}
(Table~\ref{tab:sampled}; run with and without \texttt{-{}-quadrupole});
\texttt{scripts/\allowbreak sweep\_\allowbreak multipole\_\allowbreak order\_\allowbreak tidal.\allowbreak py} (the tidal order shift) and
\texttt{scripts/\allowbreak sweep\_multipole\_order\_tidal\_camels.\allowbreak py} (the same on the real snapshot, the two
tables of Section~\ref{sec:order}, and Appendix~\ref{app:surface} with
\texttt{-{}-opening-mode surface});
\texttt{scripts/\allowbreak eval\_\allowbreak tidal\_\allowbreak eigenmetrics.\allowbreak py} (Section~\ref{sec:physics} and
Appendix~\ref{app:surface}); \texttt{scripts/\allowbreak diagnose\_\allowbreak aggregate\_\allowbreak vs\_\allowbreak percell.\allowbreak py}
(Appendix~\ref{app:percell});
\texttt{scripts/\allowbreak scaling.\allowbreak py} (Table~\ref{tab:time} and the opening-angle error curve);
\texttt{scripts/\allowbreak benchmark\_\allowbreak memory.\allowbreak py} (Tables~\ref{tab:mem} and~\ref{tab:memdeep}).
The theory note \texttt{docs/\allowbreak DERIVATION.\allowbreak md} records the bound's assumptions and closure, and
\texttt{docs/\allowbreak PROVENANCE.\allowbreak md} maps every number in this paper to the artifact that produced it,
including every defect listed above.
\end{sloppypar}

\noindent\textbf{Numerical reproducibility of the sampled statistics.} The far field is accumulated
with \texttt{index\_add\_} on the GPU, whose reduction order is not fixed, so repeated identical runs
do not return identical digits. Measured on three identical $2.3$M runs: the $p{=}0$ statistics agree
to ${\sim}10^{-7}$ relative (float32 round-off), while the $p{=}1$ statistics vary at the
${\sim}0.4\%$ level --- larger only because the error being measured is $\sim3\times$ smaller, not
because the quadrupole path is unstable. Both are well below the query-sampling uncertainty above
($1.2$--$1.7\%$), which is the spread we quote. Where this paper says a re-run ``reproduced'' an
earlier number, it means agreement at the quoted precision, not bit-identity.

\noindent\textbf{Hardware.} $1\times$NVIDIA A800 80\,GB. Every number in this paper is from a single
device; the $2.3$M row is a forward pass, not a distributed training run.

\noindent\textbf{Code version and tests.} All reported numbers were produced from repository state
tagged \texttt{v1.7-paper} (\texttt{git tag -n v1.7-paper}), which carries the D8 geometry fix and the
15-seed re-measurement of Table~\ref{tab:vjp}. \texttt{v1.6-paper} has the same geometry fix but the
previous 12-seed version of that table.
\textbf{The \texttt{v1.0-paper} and \texttt{v1.1-paper} tags predate D6 and D8 respectively and must
not be used for any number in this paper; \texttt{v1.4-paper} lacks the crossover measurement of
Section~\ref{sec:scaling} and the holdout of Appendix~\ref{app:percell}; \texttt{v1.3-paper} is
the same manuscript under the operator's former name and must not be cited by it; and
\texttt{v1.2-paper} predates the post-D8 cost
re-measurement,
the corrected real-data order-shift numbers and the real-data gradient experiment, and must not be used
for those.} Dependencies are pinned in \texttt{pyproject.toml}. The CPU
test suite passes in full ($81$ tests, including the $11$ that pin the second-derivative reference to
the autograd Jacobian); \texttt{tests/\allowbreak test\_\allowbreak camels\_\allowbreak adapter.\allowbreak py} additionally requires \texttt{h5py}
and is the one module that cannot be collected without it.

